\begin{table}[t]
\centering
\caption{Stating a pricing convention in one sentence (E10).}
\label{tab:installed}
\small
\begin{tabular}{llrrc}
\toprule
Stated convention & Ending & Blocks & Effect & 95\% CI \\
\midrule
none & .33 &  91 & -1.46 & [-3.65, 1.12] \\
whole numbers & .33 & 100 & -2.19* & [-4.10, -0.81] \\
end in .33 & .33 & 107 & 36.02* & [29.79, 42.31] \\
end in .77 & .33 & 102 & 1.65 & [-0.35, 3.55] \\
none & .77 &  91 & 2.32* & [0.22, 4.72] \\
whole numbers & .77 & 100 & -2.43* & [-4.14, -0.96] \\
end in .33 & .77 & 107 & 1.27* & [0.43, 2.37] \\
end in .77 & .77 & 102 & 33.02* & [26.84, 38.23] \\
\bottomrule
\end{tabular}
\vspace{2pt}
\parbox{\linewidth}{\footnotesize The sentence says that retail prices in a fictional country conventionally end in .33, end in .77, or are whole numbers. Effect is the ending's own effect minus the mean of the other endings in its ten-cent range, so a local price trend cannot produce it. A learned-convention account predicts a positive entry exactly where the sentence names that ending.}
\end{table}
